% TOP_PROBLEM: {"id": 5100063, "title": "Elliptic-billiard invariant k_{904,a}", "category_id": 11, "category": {"id": 11, "name": "dynamical_systems", "display_name": "Dynamical Systems", "description": "Problems about long-term behavior of deterministic systems, Hamiltonian dynamics, and periodic orbits.", "slug": "dynamical-systems", "order_index": 11, "created_at": "2026-07-31T15:26:25.671Z"}, "status": "open"}
% FIRST_POSTED: null
% ENTRY_KIND: statement_only
% Imported from: batch02.tex; UnsolvedMath ID 5100063
\documentclass[11pt]{article}
\usepackage{amsmath,amssymb,mathtools}
\usepackage{fontspec,unicode-math}
\setmainfont{Latin Modern Roman}
\setsansfont{Latin Modern Sans}
\setmathfont{Latin Modern Math}
\usepackage{xurl,hyperref}
\newcommand{\sref}[2]{\hyperlink{source:#1:#2}{[S#2]}}
\newcommand{\cataloguescope}[1]{\paragraph{Mathematical question.}#1}
\begin{document}
% UnsolvedMath ID 5100063; source catalogue code AMR-050-0063
\setcounter{section}{5100062}
\section{Elliptic-billiard invariant $k_{904,a}$}\label{5100063}
{\small\sffamily UnsolvedMath ID 5100063 · Dynamical Systems}
\subsection{Definitions and mathematical statement}

\paragraph{Shared notation.}$\mathbb N=\{1,2,\ldots\}$, and $\mathbb Z,\mathbb Q,\mathbb R,\mathbb C$ denote the integers, rationals, reals and complex numbers. Any other conventions are specified locally.


Fix $a>b>0$ and the ellipse $E:x^2/a^2+y^2/b^2=1$, with foci $f_1,f_2$. Fix an inner confocal elliptic caustic $C$ admitting a Poncelet family of $N$-periodic billiard polygons $P=(P_i)_{i=1}^N$: vertices lie on $E$, sides are tangent to $C$, and consecutive sides satisfy the billiard reflection law. Indices are cyclic and polygons retain their cyclic order. The family and its caustic stay fixed while the starting vertex varies. All areas are signed: for $W=(W_i)$,
\[\mathcal A(W)=\frac12\sum_{i=1}^N(x_i y_{i+1}-x_{i+1}y_i),\qquad W_i=(x_i,y_i).\]
Write $A=\mathcal A(P)$. Unit-circle inversion about $F$ is $I_F(X)=F+(X-F)/\lVert X-F\rVert^2$; set $P_j^\dagger=(I_{f_j}(P_i))$ and $A_j^\dagger=\mathcal A(P_j^\dagger)$, for $j=1,2$. Quotients are considered where their denominators and the polygon constructions are defined. Let $P'_i$ be the intersection of the tangents to $E$ at $P_i$ and $P_{i+1}$, and let $P''_i$ be the contact point of $P_iP_{i+1}$ with $C$. Their polygons are $P',P''$, with signed areas $A',A''$. Write ${A'_j}^\dagger=\mathcal A(I_{f_j}(P'))$ and ${A''_j}^\dagger=\mathcal A(I_{f_j}(P''))$.
\paragraph{Statement recorded in the supplied source.}
Prove that, for odd $N$, the quantity ${A'_1}^\dagger{A'_2}^\dagger$ is constant as $P$ varies through the family. \sref{5100063}{2}
\subsection{Short English statement}
Show that, for odd $N$, the quantity ${A'_1}^\dagger{A'_2}^\dagger$ is constant as $P$ varies through the family. The constant may depend on the fixed ellipse and caustic. These are the experimentally recorded assertions of the cited 2020 version.
\subsection{Sources}
\begin{description}
\item[\hypertarget{source:5100063:1}{S1}] ULAM.AI and the UnsolvedMath Contributors, UnsolvedMath: A Curated Collection of Open Mathematics Problems, record 5100063 (AMR-050-0063). CC BY 4.0. \url{https://huggingface.co/datasets/ulamai/UnsolvedMath}.
\item[\hypertarget{source:5100063:2}{S2}] Dan Reznik, Ronaldo Garcia and Jair Koiller, \emph{Eighty New Invariants in the Elliptic Billiard}, arXiv:2004.12497v11 (29 October 2020), Section 2 and Section 3.9, Table 10, row $k_{904,a}$, pp. 3, 9--12. \url{https://arxiv.org/abs/2004.12497v11}.
\end{description}
\label{end:5100063}
\end{document}
